\newpage
\phantomsection
\label{prf:higher-diff-two-n}
\LRAProofFor{prop:higher-diff-two-n}

\begin{remark*}[Return]
\hyperref[prop:higher-diff-two-n]{Return to Theorem}
\end{remark*}

\begin{proposition*}[Higher differences of \(2^n\)]
All higher forward differences of \(2^n\) are again \(2^n\).
\end{proposition*}

\begin{proof}[Professional Standard Proof]
\LRAProofBodyStart
By induction. The base case \(k = 0\) is \(\Delta^0(2^n) = 2^n\) by
definition. For the inductive step, assume \(\Delta^k(2^n) = 2^n\). Then
\[
\Delta^{k+1}(2^n) = \Delta(\Delta^k(2^n)) = \Delta(2^n) = 2^n. \qedhere
\]
\end{proof}

\begin{proof}[Detailed Learning Proof]
\LRAProofBodyStart
By induction. The base case \(k = 0\) is \(\Delta^0(2^n) = 2^n\) by
definition. For the inductive step, assume \(\Delta^k(2^n) = 2^n\). Then
\[
\Delta^{k+1}(2^n) = \Delta(\Delta^k(2^n)) = \Delta(2^n) = 2^n. \qedhere
\]
\end{proof}

\begin{remark*}[Proof structure]
\end{remark*}

\begin{dependencies}
\begin{itemize}
  \item \hyperref[def:higher-differences]{Higher Differences}.
  \item \hyperref[prop:discrete-exponential-rule]{Discrete Exponential Rule}.
\end{itemize}
\end{dependencies}
\clearpage
